\documentclass[tikz,border=8pt]{standalone}
\usepackage{tikz}
\usepackage{amsmath}
\usepackage{amssymb}
\usepackage{pgfplots}
\pgfplotsset{compat=1.18}
\usepackage{ctex}
\usetikzlibrary{calc,positioning,arrows.meta}

% LC 26 Remove Duplicates from Sorted Array
% 双指针演化：4个快照纵向排列

\begin{document}
\begin{tikzpicture}

%% ── 公共参数 ──────────────────────────────────────────────
\def\cellsize{0.8}       % 格子边长
\def\gap{3.6}            % 快照纵向间距

%% ── 数组数据：[1,1,2,2,3]，共5格 ─────────────────────────
% 数组下标 0..4，x 偏移 = index * cellsize

%% ─────────────────────────────────────────────────────────
%% 快照辅助命令（每行调用）
%  #1 = yshift（快照基准 y）
%  #2 = r 指针位置 (0-based index)
%  #3 = w 指针位置 (0-based index)
%  #4 = 数组内容（5个 token，空格分隔）
%  #5 = 说明文字

% 为避免 \newcommand 复杂性，手动写4个快照

%% ── 颜色说明 ──
% w位置格子：red!15   r位置格子：blue!15   其余：white

%% ─────────────────────────────────────────────────────────
%% 标题
\node[font=\large\bfseries] at (3.2, 1.0)
  {LC 26：双指针去重（Remove Duplicates）};

%% ─────────────────────────────────────────────────────────
%% 辅助：画一个格子
% \drawcell{x}{y}{value}{fillcolor}
\newcommand{\drawcell}[4]{%
  \node[draw, minimum size=\cellsize cm, inner sep=0pt,
        fill=#4, font=\small\bfseries] at (#1,#2) {#3};
}

%% ── 快照 0：初始状态，r=0, w=0 ──────────────────────────
\begin{scope}[yshift=0cm]
  % 格子 (values: 1 1 2 2 3)
  \drawcell{0}{0}{1}{blue!15}   % index 0: r=0, w=0 双指针重叠
  \drawcell{0.8}{0}{1}{white}
  \drawcell{1.6}{0}{2}{white}
  \drawcell{2.4}{0}{2}{white}
  \drawcell{3.2}{0}{3}{white}
  % 下标
  \foreach \i in {0,1,2,3,4}{
    \node[font=\tiny, gray] at (\i*0.8, -0.55) {\i};
  }
  % r 指针（蓝色↑）
  \draw[->,>=Stealth,blue!80!black,thick]
    (0, -0.9) -- (0, -0.5);
  \node[font=\scriptsize,blue!80!black,below] at (0,-0.95) {r=0};
  % w 指针（红色↑）—— 与 r 重叠，稍微左移显示
  \draw[->,>=Stealth,red!80!black,thick]
    (0.15, -1.3) -- (0.15, -0.5);
  \node[font=\scriptsize,red!80!black,below] at (0.15,-1.35) {w=0};
  % 说明
  \node[draw=gray!50,fill=white,thin,rounded corners=2pt,
        font=\scriptsize,inner sep=3pt,align=left]
    at (5.6, -0.4) {初始：r=0, w=0};
\end{scope}

%% ── 快照 1：r=1 与 arr[w] 相同，r 前进，w 不动 ─────────
\begin{scope}[yshift=-2.8cm]
  \drawcell{0}{0}{1}{red!15}    % w=0
  \drawcell{0.8}{0}{1}{blue!15} % r=1
  \drawcell{1.6}{0}{2}{white}
  \drawcell{2.4}{0}{2}{white}
  \drawcell{3.2}{0}{3}{white}
  \foreach \i in {0,1,2,3,4}{
    \node[font=\tiny, gray] at (\i*0.8, -0.55) {\i};
  }
  \draw[->,>=Stealth,blue!80!black,thick]
    (0.8, -0.9) -- (0.8, -0.5);
  \node[font=\scriptsize,blue!80!black,below] at (0.8,-0.95) {r=1};
  \draw[->,>=Stealth,red!80!black,thick]
    (0, -1.3) -- (0, -0.5);
  \node[font=\scriptsize,red!80!black,below] at (0,-1.35) {w=0};
  \node[draw=gray!50,fill=white,thin,rounded corners=2pt,
        font=\scriptsize,inner sep=3pt,align=left]
    at (5.6, -0.4) {arr[r]=arr[w]=1，重复\\r 前进，w 不动};
\end{scope}

%% ── 快照 2：r=2 与 arr[w] 不同，写入 w+1，w 前进 ───────
\begin{scope}[yshift=-5.6cm]
  \drawcell{0}{0}{1}{white}
  \drawcell{0.8}{0}{2}{red!15}  % w=1 (已写入)
  \drawcell{1.6}{0}{2}{blue!15} % r=2
  \drawcell{2.4}{0}{2}{white}
  \drawcell{3.2}{0}{3}{white}
  \foreach \i in {0,1,2,3,4}{
    \node[font=\tiny, gray] at (\i*0.8, -0.55) {\i};
  }
  \draw[->,>=Stealth,blue!80!black,thick]
    (1.6, -0.9) -- (1.6, -0.5);
  \node[font=\scriptsize,blue!80!black,below] at (1.6,-0.95) {r=2};
  \draw[->,>=Stealth,red!80!black,thick]
    (0.8, -1.3) -- (0.8, -0.5);
  \node[font=\scriptsize,red!80!black,below] at (0.8,-1.35) {w=1};
  \node[draw=gray!50,fill=white,thin,rounded corners=2pt,
        font=\scriptsize,inner sep=3pt,align=left]
    at (5.6, -0.4) {arr[r]=2 $\neq$ arr[w]=1\\写 arr[w+1]=2，w++};
\end{scope}

%% ── 快照 3：最终状态，r=4, w=2，结果长度 = w+1 = 3 ──────
\begin{scope}[yshift=-8.4cm]
  \drawcell{0}{0}{1}{white}
  \drawcell{0.8}{0}{2}{white}
  \drawcell{1.6}{0}{3}{red!15}  % w=2 (最终写入位置)
  \drawcell{2.4}{0}{2}{gray!20}
  \drawcell{3.2}{0}{3}{blue!15} % r=4 (终点)
  \foreach \i in {0,1,2,3,4}{
    \node[font=\tiny, gray] at (\i*0.8, -0.55) {\i};
  }
  \draw[->,>=Stealth,blue!80!black,thick]
    (3.2, -0.9) -- (3.2, -0.5);
  \node[font=\scriptsize,blue!80!black,below] at (3.2,-0.95) {r=4};
  \draw[->,>=Stealth,red!80!black,thick]
    (1.6, -1.3) -- (1.6, -0.5);
  \node[font=\scriptsize,red!80!black,below] at (1.6,-1.35) {w=2};
  % 结果框
  \node[draw=gray!50,fill=white,thin,rounded corners=2pt,
        font=\scriptsize,inner sep=3pt,align=left]
    at (5.6, -0.4) {结束：w=2，返回 w+1=\textbf{3}\\前3格 [1,2,3] 即去重结果};
  % 框选前3格
  \draw[green!60!black,thick,rounded corners=2pt]
    (-0.42, -0.42) rectangle (2.02, 0.42);
  \node[font=\tiny,green!60!black] at (0.8, 0.6) {去重结果};
\end{scope}

%% ── 图例 ─────────────────────────────────────────────────
\begin{scope}[yshift=-11.0cm]
  \node[draw,minimum size=0.5cm,fill=blue!15,inner sep=0pt]  at (0.0,0) {};
  \node[font=\scriptsize,right] at (0.32,0) {read 指针 r（蓝）};
  \node[draw,minimum size=0.5cm,fill=red!15,inner sep=0pt]   at (3.0,0) {};
  \node[font=\scriptsize,right] at (3.32,0) {write 指针 w（红）};
\end{scope}

\end{tikzpicture}
\end{document}
